\documentclass[../../../include/open-logic-section]{subfiles}

\begin{document}

\olfileid{sth}{z}{infinity-again}	
\olsection{نامتناهی}

اکنون اصول موضوع کافی در اختیار داریم تا تضمین کنیم بی‌نهایت مجموعه وجود
دارد (اگر اصلاً مجموعه‌ای وجود داشته باشد). زیرا فرض کنید مجموعه‌ای وجود
دارد و در نتیجه~$\emptyset$ وجود دارد (بنا
بر~\olref[sth][z][sep]{prop:emptyexists}). اکنون برای هر مجموعهٔ~$x$،
مجموعهٔ~$x \cup \{x\}$ بنا
بر~\olref[sth][z][pairs]{prop:pairsconsequences} وجود دارد. پس با چند بار
اعمال این عمل، مجموعه‌های زیر را به دست می‌آوریم:
\begin{enumerate}
	\item[0.] $\emptyset$ %& $0$ & stage $0$\\
	\item[1.] $ \{\emptyset\}$ %& $1$ & stage $1$\\
	\item[2.] $\{\emptyset, \{\emptyset\}\}$ %& $2$ & stage $2$\\
	\item[3.] $\{\emptyset, \{\emptyset\}, \{\emptyset, \{\emptyset\}\}\}$ %&$3$ & stage $3$\\
	\item[4.] $\{\emptyset, \{\emptyset\}, \{\emptyset, \{\emptyset\}\}, \{\emptyset, \{\emptyset\}, \{\emptyset, \{\emptyset\}\}\}\}$% & $4$ & stage $4$
\end{enumerate}
و می‌توانیم بررسی کنیم که هریک از این مجموعه‌ها متمایز است.

شماره‌گذاری را به چند دلیل از~$0$ آغاز کرده‌ایم. یکی از آن‌ها این است که
بررسی این نکته چندان دشوار نیست: مجموعه‌ای که «$n$» نامیده‌ایم دقیقاً
$n$ عضو دارد و (به‌طور شهودی) در مرحلهٔ~$n$ تشکیل می‌شود.

اما این کار \emph{بی‌نهایت مجموعه} در اختیارمان می‌گذارد، ولی تضمین
نمی‌کند که یک \emph{مجموعهٔ نامتناهی}، یعنی مجموعه‌ای با بی‌نهایت عضو،
وجود داشته باشد. و این نکته واقعاً مهم است: تا هنگامی که مجموعه‌ای
نامتناهی به معنای ددکیند نیابیم، نمی‌توانیم جبری ددکیندی بسازیم. اما
می‌خواهیم جبری ددکیندی داشته باشیم تا آن را مجموعهٔ اعداد طبیعی تلقی کنیم.
(مقایسه کنید با~\olref[sfr][infinite][dedekindsproof]{sec}.)

نکتهٔ مهم این است که اصول موضوعی که تاکنون وضع کرده‌ایم وجود هیچ مجموعهٔ
نامتناهی‌ای را تضمین \emph{نمی‌کنند}. بنابراین باید اصل موضوع تازه‌ای وضع
کنیم:

\begin{axiom}[نامتناهی]
مجموعه‌ای مانند~$I$ وجود دارد به‌طوری‌که~$\emptyset \in I$ و هرگاه
$x \in I$، آنگاه~$x \cup \{x\} \in I$.
\begin{align*}
	\exists I( & (\exists o \in I)\forall x\ x \notin o \land {}\\
	& (\forall x \in I)(\exists s \in I)\forall z(z \in s \liff (z \in x \lor z = x)))
\end{align*}
\end{axiom}

به‌آسانی می‌توان دید مجموعهٔ~$I$ که اصل موضوع نامتناهی در اختیارمان
می‌گذارد، نامتناهی ددکیندی است. عنصر ممتاز آن~$\emptyset$ است و تزریق
روی~$I$ با~$s(x)=x\cup\{x\}$ داده می‌شود. اکنون
\olref[sfr][infinite][dedekind]{thm:DedekindInfiniteAlgebra} نشان داد
چگونه از یک مجموعهٔ نامتناهی ددکیندی، جبری ددکیندی استخراج کنیم؛ و ما آن
را مجموعهٔ اعداد طبیعی تلقی خواهیم کرد. به‌بیان دقیق‌تر:
\begin{defn}\ollabel{defnomega}
فرض کنید~$I$ هر مجموعه‌ای باشد که اصل موضوع نامتناهی در اختیارمان می‌گذارد.
تابع~$s$ را با~$s(x)=x\cup\{x\}$ تعریف کنید. بگذارید
$\omega = \closureofunder{s}{\emptyset}$. اعضای~$\omega$ را
\emph{اعداد طبیعی} می‌نامیم و می‌گوییم~$n$ حاصلِ $n$ بار اعمال~$s$
بر~$\emptyset$ است.
\end{defn}

اکنون می‌توانید به چند بند پیش بازگردید و بررسی کنید که مجموعهٔ موسوم به
«$n$» \emph{به‌عنوان} عدد~$n$ تلقی خواهد شد.

اهمیت این قرارداد را در~\olref[sth][z][nat]{sec} بررسی خواهیم کرد. فعلاً
این قرارداد ما را قادر می‌سازد نتیجه‌ای شهودی را اثبات کنیم:

\begin{prop}\ollabel{naturalnumbersarentinfinite}
هیچ عدد طبیعی‌ای نامتناهی ددکیندی نیست.
\end{prop}

\begin{proof}
اثبات با استقراست؛ یعنی با استفاده
از~\olref[sfr][infinite][induction]{thm:dedinfiniteinduction}. آشکار است
که~$0=\emptyset$ نامتناهی ددکیندی نیست. برای گام استقرا عکس نقیض را ثابت
می‌کنیم: اگر (به‌فرض محال) $s(n)$ نامتناهی ددکیندی باشد، آنگاه~$n$
نامتناهی ددکیندی است.

پس فرض کنید~$s(n)$ نامتناهی ددکیندی باشد؛ یعنی یک~!!{injection}~$f$ وجود
دارد به‌طوری‌که
$\ran{f}\subsetneq \dom{f} = s(n) = n \cup \{n\}$. دو حالت را باید
بررسی کنیم.

\emph{حالت ۱: $n \notin \ran{f}$.} پس~$\ran{f} \subseteq n$ و
$f(n) \in n$. بگذارید~$g = \funrestrictionto{f}{n}$؛ اکنون
$\ran{g} = \ran{f} \setminus \{f(n)\} \subsetneq n = \dom{g}$.
ازاین‌رو~$n$ نامتناهی ددکیندی است.

\emph{حالت ۲: $n \in \ran{f}$.} یک
$m \in \dom{f} \setminus \ran{f}$ ثابت بگیرید و تابع~$h$ با دامنهٔ
$s(n)=n\cup\{n\}$ را چنین تعریف کنید:
\[
h(x) =
	\begin{cases}
		f(x) & \text{اگر }f(x) \neq n\\
		m & \text{اگر }f(x)=n
	\end{cases}
\]
پس~$h$ و~$f$ همه‌جا یکسان‌اند، جز آنکه
$h(f^{-1}(n)) = m \neq n = f(f^{-1}(n))$. چون~$f$ !!a{injection} است،
$n \notin \ran{h}$؛ و~$\ran{h} \subsetneq \dom{h} = s(n)$. اکنون با
استدلال حالت~۱، $n$ نامتناهی ددکیندی است.
\end{proof}

بااین‌همه، این پرسش باقی می‌ماند که اصل موضوع نامتناهی را چگونه می‌توان
\emph{توجیه} کرد. پاسخ کوتاه آن است که باید اصل دیگری به روایتی بیفزاییم
که تاکنون بیان کرده‌ایم. آن اصل چنین است:
\begin{enumerate}
	\item[] \stagesinf. یک مرحلهٔ نامتناهی وجود دارد؛ یعنی مرحله‌ای وجود دارد که (الف) مرحلهٔ نخست نیست، (ب) مراحلی پیش از آن وجود دارند، ولی (ج) هیچ پیشینِ بی‌واسطه‌ای ندارد.
\end{enumerate}
اصل موضوع نامتناهی به‌طور سرراست از این اصل نتیجه می‌شود. می‌دانیم که عدد
طبیعی~$n$ در مرحلهٔ~$n$ تشکیل می‌شود. پس مجموعهٔ~$\omega$ در نخستین
مرحلهٔ نامتناهی تشکیل می‌شود. و خود~$\omega$ گواه اصل موضوع نامتناهی است.

اما این صرفاً ما را به پرسش چگونگی توجیه~\stagesinf{} بازمی‌گرداند. همانند
\stagessucc، این اصل بخشی صریح از روایتی نبود که دربارهٔ سلسله‌مراتب
انباشتی-تکراری بیان کردیم. فراتر از آن، هیچ چیز در خود اندیشهٔ
سلسله‌مراتب تکراری، که در آن مجموعه‌ها مرحله‌به‌مرحله تشکیل می‌شوند، ما را
وادار نمی‌کند فرایند را شامل مرحله‌ای \emph{نامتناهی} بدانیم. کاملاً
سازگار به نظر می‌رسد که مراحل را مانند اعداد طبیعی مرتب بدانیم.

بااین‌حال، این امر مسئله‌ای آشکار پدید می‌آورد. در
\olref[sfr][infinite][dedekindsproof]{sec}، «اثبات» ددکیند بر وجود
مجموعه‌ای نامتناهی ددکیندی (از اندیشه‌ها) را بررسی کردیم. شاید آن اثبات
چندان رضایت‌بخش به نظر نرسیده باشد. اما اگر تصور تکراریِ مجموعه (یا
«قوانین اندیشه») اصل~\stagesinf{} را «بر ما تحمیل نکند»، همچنان هیچ
توجیه درونی‌ای برای این ادعا نداریم که مجموعه‌ای نامتناهی ددکیندی وجود دارد.

البته در این‌باره مطالب بسیار بیشتری می‌توان گفت. اما امید است اکنون در
نقطه‌ای باشید که بتوانید بیندیشید توجیه یک اصل موضوع (یا اصل) چه چیزهایی
ممکن است \emph{لازم داشته باشد}. در ادامه~\stagesinf{} را صرفاً مفروض
خواهیم گرفت.

\end{document}
